\newpage
\section{Gaussian Computations}
We start by demystifying the relative density function $r(x,t)$, specifically in the case that $q$ is Gaussian.

\subsection{Centering and affine invariance}
Let $q$ be the $\pi_t$-th iterate of the algorithm, and for now we don't make any assumptions on $q$.
First, we show that the variance term, which is basically the $D_{f_\alpha}$-divergence between $\PP_{X,T_0}$ and $\PP_X \otimes \PP_{T_0}$, is affine-invariant.
Recall that $X=\lambda\sum_{k=0}^{\infty}(1-\lambda)^kT_k$, where $T_0,T_1,\ldots$ are i.i.d. samples from $q$.
Let $\psi(x)=Ax+b$ be an invertible affine map, and define
\begin{align}
    \bar T_k=\psi(T_k),
    \qquad
    \bar X=\lambda\sum_{k=0}^{\infty}(1-\lambda)^k\bar T_k .
\end{align}
Since $\lambda\sum_{k=0}^{\infty}(1-\lambda)^k=1$, we have
$
    \bar X
    =
    A X+b
    =
    \psi(X)
$.

\begin{proposition}[Affine invariance]
    \label{prop:affine_invariance_f_divergence}
    For every convex generator $f$ with $f(1)=0$,
    \begin{align}
        D_f
        \left(
        \PP_{\bar X,\bar T_0}
        \middle\|
        \PP_{\bar X}\otimes \PP_{\bar T_0}
        \right)
        =
        D_f
        \left(
        \PP_{X,T_0}
        \middle\|
        \PP_X\otimes \PP_{T_0}
        \right).
    \end{align}
\end{proposition}
\begin{proof}
    Define $\Psi(x,t)=(\psi(x),\psi(t))$. Then
    \begin{align}
        \PP_{\bar X,\bar T_0}=\Psi_{\#}\PP_{X,T_0},
        \qquad
        \PP_{\bar X}\otimes \PP_{\bar T_0}=\Psi_{\#}\PP_X \otimes \PP_{T_0} .
    \end{align}
    Since $\Psi$ is a measurable bijection, then
    \begin{align}
        \frac{d\Psi_{\#}\PP_{X,T_0}}{d\Psi_{\#}\PP_X \otimes \PP_{T_0}}(\Psi(x,t))
        =
        \frac{d\PP_{X,T_0}}{d\PP_X \otimes \PP_{T_0}}(x,t)
         & =
        R\,.\qedhere
    \end{align}
\end{proof}
Consequently, whenever $q$ has mean $m$ and covariance $\Sigma\succ0$, we may apply
$\psi(x)=\Sigma^{-1/2}(x-m)$ and assume without loss of generality that it is isotropic and centered, i.e. that
\begin{align}
    \E[T_0]=0,
    \qquad
    \E[T_0T_0^\top]=I_d .
\end{align}
In particular, in the Gaussian case it is enough to study $q=\cN(0,I_d)$.


\subsection{The shape of Gaussian $r$}
In the Gaussian case, we have $q=\cN(0,I_d)$, and the next proposition gives an explicit formula for $r(x,t)$.

\begin{proposition}[Gaussian relative density]
    \label{prop:gaussian_relative_density_shape}
    Assume $q=\cN(0,I_d)$, then
    $X\sim \cN\left(0,\sigma_\lambda^2 I_d\right)$, where $\sigma_\lambda^2=\frac{\lambda}{2-\lambda}$.
    Moreover,
    \begin{align}
        r(x,t)
        =
        (1-\lambda)^{-d}
        \exp\left(
        \frac{\|x\|^2}{2\sigma_\lambda^2}
        -
        \frac{\|x-\lambda t\|^2}{2(1-\lambda)^2\sigma_\lambda^2}
        \right).
        \label{eq:gaussian_r_explicit}
    \end{align}
    Equivalently, for fixed $x$, the map $t\mapsto r(x,t)$ is proportional to a Gaussian density:
    \begin{align}
        r(x,t)
        =
        (1-\lambda)^{-d}
        \exp\left(
        \frac{2-\lambda}{2\lambda}\|x\|^2
        \right)
        \exp\left(
        -
        \frac{\lambda(2-\lambda)}{2(1-\lambda)^2}
        \left\|t-\frac{x}{\lambda}\right\|^2
        \right).
        \label{eq:gaussian_r_as_bump}
    \end{align}
\end{proposition}

\begin{proof}
    Since the $T_k$'s are independent standard Gaussians, $X$ is Gaussian with mean zero and covariance
    \begin{align}
        \Var(X)
         & =
        \lambda^2
        \sum_{k=0}^{\infty}
        (1-\lambda)^{2k} I_d
        =
        \frac{\lambda^2}{1-(1-\lambda)^2} I_d
        =
        \frac{\lambda}{2-\lambda} I_d .
    \end{align}
    Also,
    \begin{align}
        X=\lambda T_0+(1-\lambda)X',
        \qquad
        X'\sim \cN(0,\sigma_\lambda^2 I_d),
        \qquad
        X'\perp T_0 .
    \end{align}
    Hence,
    $X\mid (T_0=t)
        \sim
        \cN\left(
        \lambda t,
        (1-\lambda)^2\sigma_\lambda^2 I_d
        \right)$,
    so
    \begin{align}
        r(x,t)
        =
        \frac{p_{X\mid T_0=t}(x)}{p_X(x)}
        =
        (1-\lambda)^{-d}
        \exp\left(
        \frac{\|x\|^2}{2\sigma_\lambda^2}
        -
        \frac{\|x-\lambda t\|^2}{2(1-\lambda)^2\sigma_\lambda^2}
        \right),
    \end{align}
    which proves \eqref{eq:gaussian_r_explicit}. Completing the square in $t$ gives
    \eqref{eq:gaussian_r_as_bump}.
\end{proof}

It is also useful to record the posterior form:
\begin{align}
    r(x,t)q(t)
    =
    p_{T_0\mid X=x}(t).
\end{align}
Since $(T_0,X)$ is jointly Gaussian and
\begin{align}
    Cov(T_0,X)=\lambda I_d,
    \qquad
    Var(X)=\frac{\lambda}{2-\lambda}I_d,
\end{align}
we have
\begin{align}
    T_0\mid (X=x)
    \sim
    \cN\left(
    (2-\lambda)x,
    (1-\lambda)^2 I_d
    \right).
    \label{eq:gaussian_posterior_T_given_X}
\end{align}
Thus $r(x,\cdot)$ itself is a Gaussian-shaped bump centered at $x/\lambda$, while
$r(x,\cdot)q(\cdot)$ is the actual posterior density centered at $(2-\lambda)x$.